%!TEX root = ../thesis.tex
\setlength{\parskip}{0pt}
\chapter{基于动态哈希的电商智能检索系统的设计与实现}
\label{chap:app}

\section{引言}
前两章分别提出了面向单模态动态检索的半监督概念保留哈希方法（SCPH）与面向跨模态动态检索的多表增量哈希方法（MIH），并通过在CIFAR、MIRFlickr、MSCOCO等标准学术数据集上的实验验证了其在概念漂移场景下的有效性与稳定性。然而，学术研究的最终目的在于解决现实世界中的具体工程问题，理论方法必须在实际软件系统中落地才能充分体现其工程价值。本章旨在将SCPH与MIH算法集成到实际应用场景中，设计并实现一套完整的电商智能检索与推荐系统，验证上述方法在实际工程环境中的可用性与实用性。

在众多多媒体应用场景中，电商平台是一个典型且具有极大商业价值的动态数据环境。电商商品库规模庞大且呈现出强烈的流式增长特征：每天均有海量新品上架，商品类别不断扩展（如新品牌、新子品类、新营销活动的出现），且商品图像的拍摄风格、展示视角以及描述文本的表达方式随季节更替、流行趋势变化及营销策略调整而发生显著演化。这种典型的“概念漂移”现象会破坏传统静态检索模型所依赖的“训练数据与查询数据同分布”假设，使得基于离线全量训练的索引模型在长期运行中检索精度持续下降。若采用全量重训练以适配新数据，则需要对亿级商品库进行重新特征提取与索引重建，这不仅耗费巨大的计算与存储资源，也难以及时满足电商系统对实时性与在线更新的严格要求。

基于此，本章遵循软件工程设计规范，从系统需求分析出发，逐步完成架构设计、核心模块实现与系统测试方案的规划。该系统将SCPH算法应用于商品图像的单模态检索，支持相似商品推荐与同款发现；将MIH算法应用于商品图文对的跨模态检索，支持用户通过“以图搜商品”（拍照购）或“以文搜商品”（自然语言复杂搜索）进行灵活查询。下文将详细阐述各设计阶段的关键设计与实现细节。

\section{系统需求分析}

\subsection{业务背景}
在电商场景中，用户获取商品信息的方式正从传统的关键词检索向多模态融合方向演进。一方面，“以图搜图”已成为主流交互方式，用户在浏览某件商品时希望系统能够快速推荐款式、颜色或纹理相似的替代商品，平台方也需要在海量商品库中进行同款发现以支持比价或防止重复铺货。另一方面，“拍照购”与“复杂语义搜索”的需求日益增长：用户从社交媒体截取的穿搭图片、自然场景下拍摄的衣物照片均可作为查询输入；用户也可能输入诸如“适合秋冬季节搭配大衣的复古风粗跟马丁靴”之类的冗长自然语言描述，期望系统跨模态召回视觉上符合该语义的商品。

然而，电商数据环境具有显著的动态性与概念漂移特征。商品品类不断扩张，新款式、新子类频繁出现；季节更替导致“保暖”“轻薄”等关键词对应的商品图像分布发生剧烈变化；流行趋势的快速迭代使得某一时间段内爆火的风格（如“多巴胺穿搭”）在后续周期可能被新的潮流取代。现有检索系统在面对上述挑战时，往往难以在长期运行中保持稳定的检索精度与良好的用户体验。

\subsection{功能需求}
针对上述业务背景，本系统需提供以下核心功能。

\textbf{单模态视觉相似检索}：支持“以图搜图”场景，即用户输入商品图像后，系统能够快速从海量商品库中召回视觉风格、款式、纹理相似的商品。该功能主要服务于相似商品推荐、同款发现以及商品去重等业务。系统需能够持续适配新上架商品的类别扩展，并在仅有部分商品具备人工标注的条件下保持检索精度。

\textbf{跨模态灵活检索}：支持“以图搜文”和“以文搜图”两种跨模态查询模式。在“以图搜商品”场景下，用户上传自然场景下的衣物图片或截屏图片，系统需在电商库中找到视觉与语义均匹配的商品图文信息；在“以文搜商品”场景下，用户输入复杂的自然语言描述，系统需理解其深层语义并跨模态召回符合该描述的商品图像。两类查询均需在共享的汉明空间中完成相似度度量，且能够抵御数据分布随季节与流行趋势演化而产生的概念漂移。

\textbf{在线增量学习与索引更新}：系统需支持流式商品数据的实时接入。当新商品批次累积至设定阈值时，系统能够在不影响前台检索服务的前提下，利用增量数据在线更新SCPH与MIH模型，并为新商品生成二值哈希码写入索引。更新过程不得触发对历史全量商品哈希码的重算，以保证系统在高频上新场景下的可用性与稳定性。

\subsection{非功能需求}
在非功能层面，系统需满足以下约束。

\textbf{低延迟响应}。电商用户对检索响应时间极其敏感，在海量候选库中的召回耗时必须控制在毫秒级别。系统需充分利用哈希检索在汉明距离计算上的高效性，结合位运算实现快速相似度排序。

\textbf{高可用与可扩展性}。系统架构需支持微服务化部署与水平扩展，能够应对大促期间（如双十一、618）的高并发查询流量。各核心模块需具备独立的扩缩容能力，避免单点瓶颈。

\textbf{免重构更新能力}。在模型适应新类别或分布漂移时，必须保证历史商品哈希码的持续有效性，避免底库全量重编码。这是SCPH与MIH算法在工程化落地时的核心优势，需在系统设计中予以充分体现。

\section{系统架构设计}

\subsection{总体架构}
为满足上述需求，本系统采用微服务分层架构，自底向上分为数据接入层、特征工程层、核心哈希引擎层、索引存储层与业务服务层五大部分。各层之间通过标准化的消息队列或RPC接口通信，便于独立开发、测试与部署。系统整体架构如图\ref{fig:ecommerce-system-arch}所示。

\begin{figure}[htbp]
  \centering
  \includegraphics[width=0.9\textwidth]{app/ecommerce_system_arch}
  \caption{电商智能检索系统总体架构图}
  \label{fig:ecommerce-system-arch}
\end{figure}

\subsection{各层职责与设计}

\textbf{数据接入层}。该层负责对接电商平台的商品管理系统（Product Information Management, PIM）与用户行为日志系统。从PIM获取的数据包括新上架的商品图像、标题文本、描述属性、分类标签及上下架状态等；从日志系统可获取用户的检索请求、点击反馈等行为数据，用于后续的模型评估与权重更新。数据接入层对原始数据进行标准化清洗（如图像尺寸归一化、文本分词与去噪）、格式统一（如JSON或Protocol Buffers），并按时间窗口（如每小时或每千条）打包发送至下游特征工程层。为应对数据洪峰，该层采用消息队列Kafka进行缓冲与削峰。

\textbf{特征工程层}。该层将非结构化多模态数据转化为高维特征向量，为后续哈希学习提供统一输入。商品图像通常通过在大规模电商视觉数据集上预训练的深度卷积网络（如ResNet-50、Vision Transformer等）提取视觉特征，输出维度一般为512或1024维；商品标题与描述则通过预训练语言模型（如BERT、电商领域专用模型等）提取稠密语义特征。特征提取服务部署为无状态微服务，支持水平扩展以应对批量特征提取的算力需求。为保证与第三、四章实验设定的一致性，系统同时支持传统手工特征（如512维GIST）作为可选配置，便于在资源受限场景下运行。

\textbf{核心哈希引擎层}。该层为整个系统的算法核心，集成SCPH与MIH两类哈希模型。SCPH单模态引擎接收商品图像特征，依据第三章所述的概念保留机制与半监督学习框架，在线更新图像哈希映射函数，并输出64位二值哈希码；MIH跨模态引擎接收成对的商品图文特征，依据第四章所述的多表增量学习框架，训练跨模态哈希表并管理表权重更新与淘汰策略，输出图像与文本共享的64位哈希码。两类引擎在部署上可独立扩容，均支持按批次触发增量学习，且均保证历史哈希码在映射函数更新后仍可直接参与检索，无需重算。

\textbf{索引存储层}。该层采用向量数据库Milvus作为图像特征与哈希码的底层存储介质。每条索引记录包含商品ID、图像特征向量、64位二值哈希码以及必要的元数据（如类别、上架时间等）。Milvus支持对二值向量的汉明距离检索与Top-K近邻召回，可高效完成查询哈希码与底库哈希码的相似度计算与排序。为避免单集群容量瓶颈，索引可按商品类别或时间分区进行分片，支持水平扩展。为兼顾持久化与恢复能力，可结合对象存储与WAL机制实现索引数据的可靠落盘与快速恢复。

\textbf{业务服务层}。该层直接面向用户与运营平台，对外暴露RESTful形式的API。主要接口包括：相似商品推荐（传入商品ID，返回Top-K相似商品）、拍照购（传入用户上传图像，返回跨模态匹配的商品列表）、自然语言搜索（传入文本描述，返回跨模态匹配的商品列表）。业务服务层封装了底层的特征提取、哈希编码与索引查询逻辑，对上游提供统一的检索语义，并可根据业务规则对检索结果进行过滤、排序与截断。

\subsection{核心业务流程}
系统的核心业务流程分为商品入库与模型更新、在线用户检索两条主线。

\textbf{商品入库与模型更新流程}。新商品上架时，数据接入层将商品数据推入消息队列；特征工程层消费消息并提取图像与文本特征，将特征写入特征存储；核心哈希引擎层按设定策略（如时间窗口或累积条数）从特征存储拉取新批次数据，触发SCPH与MIH的增量学习。SCPH根据新批次中的标注与未标注样本构建概念参考矩阵与相似性矩阵，按式（3-8）求解新的投影矩阵$\mathbf{W}$；MIH根据新批次训练新的跨模态哈希表，并按最新数据更新各表权重，淘汰权重最低的表。学习完成后，引擎为新批次商品生成二值哈希码并批量写入索引存储层。全程无需访问历史商品的原始特征，也无需重算其哈希码。

\textbf{在线用户检索流程}。用户发起查询（图像或文本）时，请求经负载均衡进入业务服务层；服务层调用特征工程层提取查询特征，随后调用哈希引擎获取当前最新的哈希映射函数，将查询特征二值化为64位哈希码；业务服务层在Milvus中执行基于汉明距离的近邻检索，对候选集进行相似度排序，取Top-K结果并附加必要元数据后返回前端。为降低延迟，可对热门查询的哈希码或检索结果进行短时缓存。

\section{核心功能模块设计与实现}

\subsection{单模态图像检索模块（基于SCPH）}

单模态图像检索模块承担视觉层面的相似推荐与同款去重任务。该模块以流式上新的商品图像特征为输入，输出紧凑的二值哈希码及更新后的单模态哈希映射函数。模块内嵌SCPH算法引擎，重点解决品类扩张时的免重编码更新与弱监督环境下的语义保持两个工程难题。

\textbf{概念保留与免重编码更新}。在系统初始化阶段，模块根据当前已观测的类别数量$|L|$与哈希码长度$r$，按式（3-2）与（3-3）构造Hadamard矩阵并为各类别随机分配唯一的概念编码。概念编码一经分配即保持固定。当新款式或新类别的商品上架并触发模型更新时，模块在新批次数据上构建成对相似性矩阵$\mathbf{S}$与概念参考矩阵$\mathbf{C}$，按式（3-8）求解新的投影矩阵$\mathbf{W}$。由于概念编码在汉明空间中为各类别划定了固定区域，新学习的$\mathbf{W}$仅负责将新批次样本投影至对应区域，而不改变已有类别的投影关系。因此，历史商品的哈希码在$\mathbf{W}$更新后依然有效，彻底免除了全库重编码的开销。这一机制在电商高频上新场景下尤为重要，可保障检索服务在高负载下的高可用性。

\textbf{半监督机制与运营成本降低}。电商平台每天有大量新品涌入，且中小商家的商品往往缺乏精细、准确的人工标注。若采用全监督动态哈希，在标注稀缺时模型难以学习有效的语义结构。本模块实现了SCPH的KNN伪标签推断子模块：对于新批次中的未标注样本，模块在特征空间中计算其与已标注样本的欧氏距离（式3-4），选取K近邻并采用多数投票（式3-5）为其分配伪标签；随后结合真实标签与伪标签构建概念参考矩阵$\mathbf{C}$，参与目标函数（式3-6）的优化。该机制使得系统在仅有少量头部商家提供高质量标注的条件下，仍能充分利用海量未标注商品进行稳健的哈希学习，显著降低了平台对人工运营的依赖。

\subsection{跨模态图文检索模块（基于MIH）}

跨模态图文检索模块旨在支撑“以图搜商品”和“以文搜商品”两类业务。该模块内嵌MIH算法，需联合学习图文共享的哈希空间，并抵御数据分布随季节、流行趋势演化而产生的概念漂移与灾难性遗忘。

\textbf{多表并行架构与灾难性遗忘缓解}。电商数据分布具有强烈的周期性漂移：冬季“保暖”对应厚重羽绒服，春季“保暖”则可能对应轻薄防风外套；某年爆火的“多巴胺穿搭”在次年可能被新风格取代。若采用单一跨模态哈希模型的在线更新，模型在学习新趋势后易覆盖历史语义结构，导致用户在反季搜索或检索历史潮流商品时无法获得准确结果。本模块实现了MIH的多哈希表系统：引擎在内存中同时维护$K$组（如$K=5$）独立的跨模态哈希表，每组表由图像与文本的哈希映射函数及对应权重组成；各表分别基于不同时间阶段的数据训练得到。增量更新时，新批次数据仅用于训练新表并加入系统，旧表不被覆盖；当表数量超过$K$时，按权重淘汰权重最低的表。该架构使系统能够并行保留不同时期的跨模态语义结构，从根本上缓解灾难性遗忘。

\textbf{双重动态加权与趋势适配}。在存在多张历史哈希表的情况下，如何决定各表对当前查询的贡献成为关键。本模块实现了基于新数据的表权重更新与查询自适应加权两种机制。表权重更新方面：当具有当前流行趋势的新商品批次入库时，后台任务使用该批次的图文对对各保留哈希表进行跨模态一致性评估（式4-9、4-10），即衡量各表将图像与文本映射至共享哈希空间后与学习得到的共享哈希码的一致性；一致性越高的表被赋予更大权重。查询自适应加权方面：在每次在线检索时，系统计算查询样本与各哈希超平面的距离（式4-12），对距离大、判别能力强的比特位赋予更大权重（式4-13），在加权汉明距离计算（式4-14）中提升其贡献。双重加权机制使系统既能保留历史多表的语义知识，又能敏锐适配当前查询与流行趋势，从而在动态变化的电商搜索中提供精准的召回排序。

\section{系统的实现与测试}

本节对系统的实现环境、关键接口设计与测试方案进行说明，为后续的完整实现与实验验证提供参考框架。

\textbf{实现环境}。系统采用Python作为主要开发语言。SCPH与MIH的核心算法逻辑基于NumPy、SciPy实现矩阵运算与闭式解求解；特征提取模块基于PyTorch或TensorFlow加载在大规模电商视觉与文本数据上预训练的ResNet、BERT等模型，支持GPU加速；业务服务层采用FastAPI或gRPC框架对外暴露RESTful或二进制RPC接口，便于前端与移动端调用；索引存储采用Milvus，通过PyMilvus客户端完成图像特征与二值哈希码的批量写入、分区管理与Top-K检索，汉明距离计算由Milvus检索引擎完成。部署方面，各层以Docker镜像形式打包，通过Kubernetes进行容器编排、服务发现与自动扩缩容，保证大促期间的高可用性。

\textbf{关键接口设计}。核心哈希引擎层对外提供两类主要接口：一是“编码接口”，输入为特征向量或特征向量列表，输出为64位二值哈希码；二是“更新接口”，输入为新批次的特征矩阵与标签信息，触发增量学习并返回更新后的模型参数。业务服务层对外提供三类检索接口：相似推荐（输入商品ID，内部获取该商品特征并编码后检索）、拍照购（输入用户上传图像，经特征提取与跨模态编码后检索）、自然语言搜索（输入文本，经特征提取与跨模态编码后检索）。各接口均支持Top-K参数与过滤条件（如类别、价格区间），便于与上游业务系统集成。

\textbf{测试方案}。系统测试从单元测试、集成测试与性能测试三个层面展开。单元测试针对SCPH的概念编码生成、伪标签推断、投影矩阵求解以及MIH的表权重计算、查询自适应权重计算等核心函数，使用小规模人工构造数据或从标准数据集采样的子集验证其正确性与数值稳定性。集成测试验证数据接入层到业务服务层的端到端数据流：模拟新商品上架触发特征提取与哈希编码写入，模拟用户查询触发编码与检索返回，检查各环节的数据格式、异常处理与边界条件。性能测试在构造的百万级商品库与模拟的高并发查询负载下，评估单次检索的P99延迟、系统吞吐量（QPS）以及模型增量更新的耗时，验证是否满足毫秒级响应与免重构更新的非功能约束。

\section{本章小结}

本章从软件工程角度设计并阐述了基于SCPH与MIH算法的电商智能检索系统。通过需求分析，明确了电商场景下全量重构成本高昂、标注稀缺与灾难性遗忘三类核心痛点，以及单模态与跨模态检索、在线增量更新等功能需求与低延迟、高可用、免重构等非功能需求。系统采用数据接入、特征工程、核心哈希引擎、索引存储与业务服务五层微服务架构，明确了各层职责与核心业务流程。

在核心功能模块方面，单模态图像检索模块依托SCPH的概念保留机制与KNN伪标签推断，实现了上新扩张时的哈希码免重构更新，并在弱监督环境下降低了运营标注成本；跨模态图文检索模块通过MIH的多表并行架构与双重动态加权机制，有效缓解了电商数据周期性演化带来的灾难性遗忘，提升了“拍照购”与“自然语言复杂搜索”的长效准确性。本章的设计与分析验证了本文提出的动态哈希检索方法在应对真实工业级概念漂移问题时的有效性，为新一代多媒体智能检索系统的落地提供了具有参考价值的软件工程实践方案。
